Pessoas com esclerose lateral amiotrófica perdem progressivamente a fala e os movimentos, mas em geral preservam o controle dos olhos, mas os dispositivos comerciais de comunicação pelo olhar têm custo fora do alcance da maior parte das famílias brasileiras. Este trabalho apresenta um aplicativo de comunicação controlado pelo olhar para \textit{tablets} Android comuns, que usa apenas a câmera frontal do aparelho e funciona sem conexão com a internet. O rastreamento ocular combina a detecção dos cantos dos olhos pelo ML Kit, a localização da pupila por uma rede YOLO26n, uma calibração por usuário baseada em regressão polinomial e um filtro de Kalman, e converte a direção do olhar em comandos de navegação e seleção em um teclado. Para reduzir o esforço de digitação, um modelo de trigramas com \textit{Stupid Backoff}, treinado em frases de comunicação aumentativa e alternativa e tolerante a erros de digitação, sugere palavras, e a mensagem composta é falada por um modelo de síntese de voz em português executado no próprio aparelho. Um módulo de métricas registra cada sessão e calcula a velocidade, o esforço e a contribuição das sugestões. O sistema foi avaliado em três testes de uso feitos por um dos autores em um Samsung Galaxy Tab S7, com cinco repetições cada e todos os comandos emitidos pelo olhar. O desempenho melhorou com a prática, e a velocidade passou a depender sobretudo do número de comandos. Com as sugestões, a economia de comandos chegou a 79,4\% e a velocidade a cerca de 3 palavras por minuto, na mesma faixa de outros sistemas de baixo custo baseados em câmera. Como trabalhos futuros, propõem-se testes com pessoas com esclerose lateral amiotrófica, um aviso sonoro da tecla selecionada e sugestões de palavras baseadas no contexto de uso.
